\documentclass[11pt]{article}
\usepackage[T1]{fontenc}
\usepackage{lmodern,microtype}
\usepackage{amsmath,amssymb,amsthm}
\usepackage[letterpaper,margin=1in]{geometry}
\usepackage{booktabs,array,enumitem}
\usepackage[colorlinks=true,linkcolor=blue!50!black,urlcolor=blue!50!black]{hyperref}
\newcommand{\GL}{\mathrm{GL}}
\newcommand{\A}{\mathbb A}
\newcommand{\Gb}{\overline G}
\newcommand{\GDb}{\overline G_D}
\newcommand{\Zb}{\overline Z}
\newcommand{\tH}{\widetilde H}
\newcommand{\tZ}{\widetilde Z}
\newcommand{\Nrd}{\operatorname{Nrd}}
\theoremstyle{definition}
\newtheorem*{prop*}{Proposed Corollary}

\title{Referee report on \texttt{mjl\_part4\_draft1.tex}\\
\normalsize ``Toward a metaplectic Jacquet--Langlands transfer, IV''}
\author{Claude (Anthropic), for David Joyner}
\date{November 7, 2026}

\begin{document}
\maketitle

\section*{0. Scope and what I read}
\paragraph{Read in full.} Part~IV, Draft~1 (669 lines), and its log;
Part~III, Draft~6 (1359 lines); the Part~III log (Entries 0--6); my
Draft~5 report on Part~III; the Part~II log; and the revision log for
Part~I Draft~6 and [JC] revised12.

\paragraph{Read in part.}
\begin{itemize}
\item \emph{Part~I, Draft~6:} lines 1--949, i.e.\ \S\S1--8 (everything
except the glossary, literature section and Appendices A, B).
\item \emph{Part~II, Draft~6:} the introduction and dependency diagram
(lines 1--140), Definition~2.1 and Remark~2.2 (measures), \S7 (``What is
left for Part~III''), and the acknowledgments. I did not check
Part~II's proofs here.
\item \emph{[FK] (scan) and Zou (arXiv v3):} I extracted the text and read
only the passages listed in \S3 below.
\end{itemize}

\paragraph{Not read.} [JC] (\texttt{metaplectic\_kazhdan\_revised12.tex}):
I compiled it but did not read it, since Part~IV does not cite it.

\paragraph{Compilation.} All of Parts~I--IV compile. In my build the page
counts are 48, 32, 23 and 11. The theorem numbers I quote below come
from the compiled \texttt{.aux} files. Part~IV's numbering is as follows (equations are numbered consecutively in all four papers):
\begin{itemize}
\item Proposition~2.1 (lattice criterion);
\item Remark~2.2;
\item Proposition~2.3 (finite-level test);
\item Remarks~2.4, 2.5;
\item Lemma~3.1 (radical);
\item Theorem~3.2 (split diagonal good elements);
\item Remarks~3.3, 3.4;
\item Definition~3.5 (nonsplit-torus problem);
\item Proposition~4.1 (Hasse principle);
\item Remark~4.2;
\item Definition~5.1 (geometric comparison);
\item Remark~5.2;
\item Proposition~6.1 (spectral).
\end{itemize}

\section*{Summary verdict}
I found no false mathematical identity in Part~IV. Every displayed
formula I checked is correct. The draft's main problem is that it does
not take account of what Part~I and [FK] already contain.
\begin{enumerate}[label=(\arabic*)]
\item \emph{The split-place theorem and the ``open'' nonsplit problem are
already in [FK].} Theorem~3.2 is [FK, \S3, Proposition] specialized to
the diagonal torus. That Proposition treats \emph{every} regular torus
$\bigoplus_iF_i^\times$, at every finite place, tame or wild. So
Definition~3.5 is not an open problem. Part~I already uses the
Proposition: in the proof of Theorem~5.6, and through its commutator
formula in Proposition~3.1.
\item \emph{The global-comparison sections ignore Parts~I and~II.}
\begin{itemize}
\item Part~I already contains a proved trace-formula comparison:
Theorem~6.6, with rational-orbit accounting (Theorem~5.6,
Corollary~5.7), measures (\S5.3), and identity terms (Proposition~6.1).
Part~II extends it (Theorem~6.2).
\item Part~IV's introduction, status table, \S4 and \S5 describe all of
this as open.
\item Definition~5.1 does not say \emph{which} two groups are being
compared. Part~I compares $\GDb$ with $\Gb$. The target identity
(31) of Part~III compares $\GDb$ with $D^\times$.
\end{itemize}
\item \emph{The draft misses a new global result that it is close to.}
With (1), Proposition~4.1 \emph{already} identifies the globally good
regular elliptic classes. Remark~4.2 says this is not yet possible.
I propose a short replacement corollary in \S1.4.
\item \emph{Section~2 adds little to Part~III.} Propositions~2.1, 2.3 and~6.1
restate Part~III, Corollaries~5.1 and~5.2. The ``inverse'' conditions
in them are redundant, because Part~III, Theorem~4.3, proves that $L_D$
is a unitary \emph{isomorphism}.
\end{enumerate}
Once these points are dealt with, the paper's genuinely new content is
small: the corollary in \S1.4, the finite-level wild test, and a precise
formulation of the $\GDb$--$D^\times$ comparison. I recommend
restructuring along those lines (\S5) before any further refereeing.

\begin{center}\small
\begin{tabular}{@{}p{0.30\linewidth}p{0.17\linewidth}p{0.45\linewidth}@{}}
\toprule
Result & Mark & Main issue\\\midrule
Prop.~2.1 (lattice criterion) & correct; minor repair & Its hypotheses are Part~III theorems; the inverse clause is redundant; it is Part~III, Cor.~5.2, plus a trivial converse\\
Rem.~2.2 & misleading & For a unitary $L_D$, the inverse condition is automatic\\
Prop.~2.3 (finite level) & correct & This is Part~III, Cor.~5.1\\
Rem.~2.5 (Adams) & correct; minor repair & The relevant condition is $n$ prime to $|Q/C|$, not to $|Q|$\\
Eq.~(7), Lemma~3.1 & correct; minor repair & $c$ should be $c_v$; ``alternating'' is misnamed; $(r,n)=1$ is automatic\\
Thm.~3.2 & correct & It is [FK, \S3, Prop.] for the diagonal torus; label \textbf{[E]}\\
Def.~3.5 \textbf{[O]} & status false & Already solved by [FK, \S3, Prop.] (see \S1.3)\\
Prop.~4.1 & correct & This is Part~III, Prop.~11.3; cite the number\\
Rem.~4.2 & incorrect claim & The identification it calls missing is in [FK, \S3]; replace it by the corollary in \S1.4\\
\S5.1 (three quantities) & minor repair & $b_D$ and $c_D$ do not depend on the choice of measures (\S1.5)\\
Ledger, Def.~5.1 & serious gap (formulation) & The groups compared are unstated; several rows are settled in Part~I\\
Prop.~6.1 & correct & This is Part~III, Cor.~5.2 again; the inverse clause is redundant\\
\bottomrule
\end{tabular}
\end{center}

\section*{1. Correctness}

\subsection*{1.1 Section 2: the character lattice}
\paragraph{Proposition~2.1: correct, minor repair.}
\begin{itemize}
\item \emph{The hypotheses are not hypotheses.} ``Unitarity, smoothness, and
finite character expansions'' are proved in Part~III: Theorem~4.3,
Lemma~4.2 and Corollary~5.1. Assuming them makes the result look
conditional. (I raised the same point about ``Hypothesis~(G)'' in my
Draft~5 report on Part~III.)
\item \emph{The inverse clause is automatic.}
\begin{itemize}
\item Since $L_D$ is a unitary isomorphism and both character families
are orthonormal bases, the coefficient of $\Theta_\rho$ in
$L_D^{-1}\Theta_\tau$ is $\overline{a_{\rho\tau}}$.
\item If every $a_{\rho\tau}\in\mathbb Z$, each row is a signed basis
vector. Distinct rows are orthogonal, so $\rho\mapsto\tau(\rho)$ is
injective.
\item Every column has norm~$1$, so every $\tau$ is hit.
\end{itemize}
Hence (iii) can be stated as ``every $a_{\rho\tau}$ is an integer''. The
sentences about the adjoint matrix should then be deleted.
\item \emph{Remark~2.2 should be rewritten.} As it stands it suggests a real
issue that unitarity already removes.
\item \emph{Status.} The implication (iii)$\Rightarrow$(i) is Part~III,
Corollary~5.2. The converse is trivial. I suggest the label
\textbf{[E]} (Part~III, Cor.~5.2) with a one-line proof of the converse,
not a new label \textbf{[R]}.
\end{itemize}
The sign argument is correct: $(L_D\Theta_\rho)(1)=c_D\dim\rho>0$ by
Part~III, Corollary~5.2.

\paragraph{Proposition~2.3 and Remark~2.4: correct.} This answers log
question~(Q3) affirmatively, and the argument is already the proof of
Part~III, Corollary~5.1. A smooth $\tilde\theta$-equivariant invariant
function is right-invariant under an open normal subgroup $N$, because
$\tH/\tZ$ is compact. Only the finitely many $\tau$ with $\tau^N\ne0$
and central character $\tilde\theta$ can occur, since $N\tZ$ has finite
index. I suggest citing Part~III, Corollary~5.1, and keeping only the
computational remark.

\paragraph{Remark~2.5 (Adams operations): correct, minor repair.} The
sentence ``when $n$ is prime to $|Q|$ \ldots\ this is the ingredient used
in the tame case of Part~III'' misdescribes Part~III, Lemma~7.5. That
lemma allows a central subgroup $C$ of order divisible by $n$, with the
central character fixed, and needs only $\gcd(n,|Q/C|)=1$. The obstruction
at a wild place is therefore $p\mid|Q/C|$. This does occur:
\begin{itemize}
\item every finite quotient of $H/Z$ maps onto a quotient of
$E^\times/F^\times E^{\times n}$ for a maximal subfield $E$;
\item that group has order $|n|_F^{-(r-1)}$ (Part~III, Lemma~3.2).
\end{itemize}
The $C_p$ example is right if read with $C=1$. It should say so.

\subsection*{1.2 Section 3: the commutator matrix}
\paragraph{Equation (7) and Lemma~3.1: correct.} I recomputed
$\sigma_c(t,u)/\sigma_c(u,t)$ directly from (5). Antisymmetry
$(a,b)(b,a)=1$ gives
\[
\prod_{i,j}(t_i,u_j)^{2c+[i\ne j]},\qquad\text{so}\qquad A=-I+(1+2c)J.
\]
This agrees with Part~I, Proposition~3.1. That proposition, following
[FK, \S3] and [KP1, Prop.~0.1.5], gives
$[x,y]=(\det x,\det y)^{1+2c_v}\prod_i(x_i,y_i)^{-1}$ on any maximal
torus. The kernel computation in Lemma~3.1 is correct. Four repairs:
\begin{enumerate}[label=(\alph*)]
\item \emph{Notation.} Write $c_v$ for the parameter. In Part~I
(Assumption~1.1) the letter $c$ denotes the integer with $n\mid r-1+2c$,
and $c_v=r'c$.
\begin{itemize}
\item Part~IV's $c$ in (5)--(6) is Part~I's $c_v$. The example
$c\equiv2\pmod3$ matches Part~I, Example~3.7, which has $c=1$ and $c_v=2$.
\item Since $2r$ is invertible modulo $n$, the congruence determines $c_v$
modulo $n$. Part~I, \S5.1, notes that $c_v$ does not depend on $v$. So
the ledger row ``dependence on $c_v$'' is vacuous.
\end{itemize}
\item \emph{Terminology.} ``Since $n$ is odd, the Hilbert symbol is
alternating: $(x,y)_n=(y,x)_n^{-1}$.'' The displayed identity is
antisymmetry, which holds for every $n$. ``Alternating'', meaning
$(x,x)=1$, is what needs $n$ odd, and it is not used.
\item \emph{A redundant hypothesis.} If $n\mid r-1+2rc$ then $\gcd(r,n)$ divides
$1$, so $(r,n)=1$ is automatic.
\item \emph{Matching Part~I's cocycle.} Log question~(Q1) asks whether
(5) agrees with Part~I. It does, and literally, not only up to a
coboundary, in the $\sigma$-model after relabeling. Part~I,
Proposition~2.3(b) and Step~1 of its proof, show that $\Gb_v$ is the
extension with cocycle $\sigma^{\,r}(\det,\det)^c$. On the diagonal
torus this is $\prod_{i<j}(a_i,a'_j)^r(\det,\det)^c$. After the
relabeling $\zeta\mapsto\zeta^{r'}$ it is exactly (5) with $c=c_v$.
Before relabeling, every commutator is raised to the power $r$, and the
radical is unchanged.
\end{enumerate}

\paragraph{Theorem~3.2: correct, but known.} The statement is
[FK, \S3, Proposition]: ``$[x,g]=1$ for all $g$ in $T$ if and only if $x$
lies in $Z_*T^n$.'' I checked this on p.~59 of the attached scan. There
$Z_*$ is the inverse image of the center of the cover, and $Z_*=Z$ by
Part~I, Lemma~2.2. The theorem should be labeled \textbf{[E]} and
attributed. The abstract's ``we compute the commutator radical'' should
become ``we recall''.

\subsection*{1.3 Definition~3.5 is not open}
[FK, \S3] states the commutator formula for the torus
$T=\bigoplus_iF_i^\times$ of an arbitrary regular element:
$[x,g]=(\det x,\det g)^{1+2c}\prod_i(x_i,g_i)_{F_i}^{-1}$. It deduces
goodness $\iff x\in Z_*T^n$ for every such torus. The proof is three
lines, and it is the proof of Part~I, Proposition~3.5(a), word for word.
\begin{itemize}
\item Write $E=\prod_iE_i$ and $N=N_{E/F}$.
\item Norm compatibility gives
$(\det x,\det g)_F=\prod_i(Nx,g_i)_{E_i}$.
\item Nondegeneracy of each $(\cdot,\cdot)_{E_i}$ (valid at every finite
place, since $\mu_n\subset E_i$) gives: $x$ is good $\iff$
$x_i\equiv(Nx)^{1+2c_v}\bmod E_i^{\times n}$ for all $i$.
\item The congruence $n\mid r-1+2rc_v$ turns this into
$x\in F^\times\prod_iE_i^{\times n}$, exactly as in Part~I.
\end{itemize}
No new ``restriction and corestriction'' analysis is needed, and the
answer is uniform in $E$. Places with $p\mid n$ are included. So
Definition~3.5 should be deleted. The paragraph after it, warning that
rational points need separate treatment, is moot: the formula is stated
on $F$-rational points.

\subsection*{1.4 Section 4: what Proposition~4.1 already gives}
\paragraph{Proposition~4.1: correct.} This is Part~III, Proposition~11.3;
I rechecked the algebra
$t/a=(a_v^k y_vN(y_v)^{-u})^n$. Log question~(Q5): the hypotheses used
are
\begin{itemize}
\item $[E:K]=r$ (in the norm trick);
\item $\gcd(n,r)=1$ (to choose $u,k$);
\item $\mu_n\subset K$ (Kummer theory over $E$).
\end{itemize}
$n$ odd is not used.

\paragraph{Remark~4.2 is wrong in its key claim.} It says the
proposition ``cannot be invoked from the good-element condition unless
that condition has first been identified with the local power condition
for the actual regular centralizer''. By \S1.3 that identification holds
at every place. I suggest replacing the remark by the following
statement, which I believe is new in the series and is the rational-orbit
accounting asked for in Part~III, \S12, obligation~(1).

\begin{prop*}[\textbf{[P]}]
Assume Part~I, Assumption~1.1. Let $G\in\{G_D,\GL(r)\}$, and let
$\gamma\in G(K)$ be regular elliptic, with $E=K(\gamma)$ a field of degree
$r$. The following are equivalent:
\begin{enumerate}[label=(\alph*)]
\item $\iota(\gamma)$ is good at every place;
\item $\iota(\gamma)$ is good at every place outside some finite set
$\Sigma$;
\item $\gamma\in K^\times E^{\times n}$.
\end{enumerate}
Moreover $s\mapsto s^n$ is a bijection from the regular elements of
$E^\times/K^\times$ onto the classes of $E^\times/K^\times$ satisfying
(a)--(c), and $W_{s^n}=W_s$ in the notation of Part~I, \S6.
\end{prop*}

\noindent\emph{Proof sketch.}
\begin{enumerate}[label=(\roman*)]
\item \emph{The local criterion at each place.} At a finite place $v$, the
centralizer of $\gamma$ is $T_v=(E\otimes_KK_v)^\times$.
\begin{itemize}
\item For $v\notin S$, combine Part~I, Proposition~2.6 (on the $D$ side),
Proposition~2.3(b), Lemma~2.2 and [FK, \S3]: goodness at $v$
$\iff\gamma\in K_v^\times T_v^n$.
\item For $v\in S$, the same holds by Part~I, Proposition~3.5(a).
\item At the complex places everything is split and $T_v^n=T_v$.
\end{itemize}
\item \emph{The equivalence.} (c)$\Rightarrow$(a)$\Rightarrow$(b) are
clear, and (b)$\Rightarrow$(c) is Proposition~4.1.
\item \emph{Injectivity of $s\mapsto s^n$.} If $s^n\in K^\times$, then
$[K(s):K]$ divides both $n$ and $r$, so it is $1$.
\item \emph{Regularity.} If $a s^n$ is regular then $s$ is not central.
\item \emph{The groups $W$.} For $\sigma\in\operatorname{Aut}(E/K)$,
$\sigma(s^n)\in K^\times s^n\iff(\sigma(s)/s)^n\in K^\times\iff
\sigma(s)\in K^\times s$, by the same injectivity. \qed
\end{enumerate}

So the geometric sides in Part~I, Propositions~6.1 and~6.2, run over
$\{[s^n]\}$, and the global ``Kummer obstruction'' disappears for the
comparison of $\GDb$ with $D^\times$.

\paragraph{\S4.1 (Kottwitz): correct.} This is consistent with
Part~III, \S10.

\subsection*{1.5 Section 5: normalizations}
\begin{itemize}
\item \emph{The formulas are right.} The formulas for $b_{D,v}$ and
$c_{D,v}$ agree with Part~III, Theorem~4.3, Corollary~5.2 and
Corollary~9.3. Log question~(Q2): yes.
\item \emph{But $b_D$ and $c_D$ do not depend on the measures.} Part~III
identifies $\tH/\tZ\cong H/Z$ canonically. The covered orbital integral
at $zt^*$ is an integral over the same quotient $H/T$ as the linear one
at $t$. So if the same Haar measure $dg$ is used on $H/Z$ and
$\tH/\tZ$, and the same $dt$ on $T/Z$:
\begin{itemize}
\item both orbital integrals in Part~III, (20), are multiplied by the
same factor, the ratio of $\mathrm{vol}_{dg}(H/Z)$ to
$\mathrm{vol}_{dt}(T/Z)$;
\item both $L^2$ norms in Theorem~4.3 are multiplied by the same constant.
\end{itemize}
Thus (20), (21) and unitarity hold with $b_D$ for \emph{every} common choice
of measures. The value $c_D$ at the identity is pointwise, so it is
measure-free. The ``third contribution'' in \S5.1 is therefore zero for
a comparison of $\GDb$ with $D^\times$. Measures matter only in the
comparison of $\GDb$ with $\Gb$, where Part~II, Definition~2.1, fixes
$\mathrm{vol}(D_v^\times/F^\times)=d(\mathrm{St})^{-1}$. Under that
choice the probability measure is $d(\mathrm{St};dg_v)\,dg_{D,v}$.
\item \emph{The product-formula remark is correct.} Here is a concrete point
for the ledger. At a place $v\notin S$ where $E\otimes K_v$ has $k_v$
factors, the map $s\mapsto s^n$ on $T_v/Z_v$ is \emph{not} injective. Its
kernel is $\mu_n^{k_v}/\mu_n$, of order $n^{k_v-1}$, even when $v\nmid n$.
\begin{itemize}
\item Contrast $v\in S$: there the power map is injective (Part~III,
Lemma~3.2).
\item This is the reason for the sum over $\{t:t^n=x^n\}$ in Zou's (3.2).
\end{itemize}
The local transfer at split places, and the bookkeeping of these kernels
against the single global $s$ of \S1.4, must be stated explicitly.
\item \emph{The ledger rows.} Several are already settled.
\begin{itemize}
\item ``Global cocycle'': Part~I, Lemma~5.1 and Proposition~5.5
(reciprocity, rational and compact splittings).
\item ``Global torus volume'': Part~I, \S5.3, uses one measure on
$\A_E^\times/\A_K^\times$ on both sides.
\item ``Local orbit'': \S1.4 above.
\item ``Singular and identity terms'': Part~I, Proposition~6.1, computes the
identity term; regular support at one place kills it; Part~II,
Theorem~5.13, treats the germ at the center.
\end{itemize}
\item \emph{Definition~5.1 must name the groups.} If the target is
Part~III, (31), i.e.\ $\GDb(\A_K)$ against $D^\times(\A_K)$, then:
\begin{itemize}
\item the $\GDb$ geometric side is Part~I, Proposition~6.1;
\item the $D^\times$ side is the classical one;
\item the orbit bijection is \S1.4;
\item at $v\in S$ the local identity is Part~III, Theorem~9.1;
\item what remains is (i) the local transfer at $v\notin S$ between
$\Gb_v$ and $\GL(r,K_v)$, which is the Flicker--Kazhdan matching, with
unit elements under $(n,N)=1$ [FK, Corollary of \S12], and (ii) the
product of all local constants.
\end{itemize}
That is a far more precise open problem than the present Definition~5.1.
\item \emph{Remark~5.2.} Add, as Part~II, \S7(iii) does, that [Li4]
assumes an invariant Paley--Wiener theorem at the archimedean places.
\end{itemize}

\subsection*{1.6 Section 6}
\paragraph{Proposition~6.1: correct.} Given integral $m_\tau$, unitarity
gives $\sum m_\tau^2=1$, and positivity at $1$ fixes the sign. This is
Part~III, Corollary~5.2, applied to one row. The hypothesis ``the inverse
coefficients are integral'' is redundant (\S1.1). The proposition is
labeled \textbf{[R]}, but no reduction is proved. It is a restatement.

\section*{2. Consistency with Parts~I--III}
\begin{enumerate}[label=(\arabic*)]
\item \emph{Description of the earlier parts.} \S1 says Part~I develops
``the specific covering cocycles, their centers, and good elements''. In
fact Part~I proves Theorems~A, B and~C:
\begin{itemize}
\item Theorem~3.9;
\item Theorem~4.5;
\item the trace-formula comparison, Theorem~6.6;
\item the packet identity, Proposition~7.6;
\item Theorem~C itself, Theorem~7.11, under $(\mathrm H_D)$.
\end{itemize}
Part~II proves Theorem~E (Theorem~5.8) and Theorem~C$'$ (Theorem~6.5,
under $(\mathrm H_D')$, Definition~6.4). The status-table row ``Global
trace-formula coefficients: \textbf{[O]}'' contradicts Part~I,
Theorem~6.6 \textbf{[P]}. It should refer specifically to the
$\GDb$--$D^\times$ comparison.
\item \emph{Cross-references.} Part~IV never cites a numbered result of
Parts~I--III. Suggested citations:
\begin{itemize}
\item $L_D$: Part~III, Def.~4.1 and Thm.~4.3.
\item Finite expansions: Part~III, Cor.~5.1.
\item The identity value and $c_D$: Part~III, Cor.~5.2, Lemma~5.4 and
Cor.~9.3.
\item The factor $b_D$ in dimensions: Part~III, Prop.~8.1.
\item The tame case: Part~III, Lemma~7.5 and Thm.~7.6.
\item The Hasse principle: Part~III, Prop.~11.3.
\item Scalar centrality: Part~I, Lemma~2.2.
\item The cocycle: Part~I, Prop.~2.3 and eq.~(1).
\item The torus commutator: Part~I, Prop.~3.1.
\item Elliptic good elements: Part~I, Prop.~3.5.
\end{itemize}
\item \emph{Central characters.} \S2.1 says ``fix $\theta$ \ldots\ and the
corresponding genuine $\tilde\theta$''. Part~III goes the other way:
it fixes $\tilde\theta$ and puts $\theta(a)=\tilde\theta(a^*)$. Since
$\theta=\tilde\theta^{\,n}$ on lifts, $\theta$ does not determine
$\tilde\theta$. Fix $\tilde\theta$ first, and write
$L_D^{(\tilde\theta)}$ where it matters.
\item \emph{Labels.} Part~I defines \textbf{[O]} as ``proof outlined'' and
\textbf{[C]} as ``conjectural''. Parts~III and IV use \textbf{[O]} for
``open problem'' and \textbf{[C]} for ``conditional''. Part~IV adds
\textbf{[R]}. One convention should be adopted for the series. I would
drop \textbf{[R]}.
\item \emph{The sign $(-1)^{r-1}$ and the link to Part~I.} Part~IV works
only on the division side and never mentions the GL-side character
relation. Your consistency report says that the equality $d_r=n$ makes
Zou's constant in (3.2) equal to $b_D$. Then Part~I's
$\Theta_{\pib_D}=(-1)^{r-1}\Theta_{\pib}\circ\jmath$ (goal (D2); Part~II,
\S7(ii)) is equivalent to integrality of $L_D$. If so, the
``wild local correspondence'' of Part~IV, \S2, \emph{is} Part~II's goal~(ii)
at wild places of $S$, and the paper should say so. I checked two
ingredients of your report against Zou:
\begin{itemize}
\item Zou, (2.2), gives the center as
$p^{-1}\{\lambda:\lambda^{2rc+r-1}\in F^{\times n}\}$, the same
convention as ours;
\item Zou, \S2.3, defines $d_{r_0}=\gcd(n,2r_0c+r_0-1)$.
\end{itemize}
So $d_r=n$ under $n\mid r-1+2rc_v$. Note, however, that a later passage
of Zou (p.~18 of arXiv v3, \S4, just after the proof of Proposition~4.2)
writes $d_r=\gcd(n,2rc-r+1)$. This appears to be a typo in Zou, but it is
worth a footnote.
\item \emph{$(n,N)=1$.} [FK, Introduction and Corollary of \S12] assume
$(n,N)=1$, with $N$ the least common multiple of the composite integers
$\le r$. Part~III, Remark~12.1, records this. Part~IV does not mention
it, although any use of Flicker--Kazhdan matching at split places
(\S1.5) needs it unless Zou or later work removes it.
\item \emph{Measures and $\Delta$.} Part~I's
$\Delta=|D(\gamma)|^{1/2}|\det\gamma|^{-(r-1)/2}$ (Definition~4.1) equals
Part~III's product (3), since $\prod_{i<j}\lambda_i\lambda_j=\det^{r-1}$.
This is consistent. For measures, see \S1.5.
\item \emph{Titles and authors.} The log says the title is ``aligned with the
rest of the series''. It is aligned with Part~III only. Parts~I and~II are
titled ``A metaplectic Jacquet--Langlands transfer for division algebras
of prime degree, I/II''. The author lines also differ:
\begin{itemize}
\item Parts~I, II: ``David Joyner and Claude'';
\item Part~III: Joyner, with assistance from Claude and ChatGPT;
\item Part~IV: Joyner, with assistance from ChatGPT only.
\end{itemize}
This is your decision. The four papers should be made consistent.
\end{enumerate}

\section*{3. Citations}
\paragraph{Checked against the attached sources.}
\begin{itemize}
\item \emph{[FK86].} Publ.\ Math.\ IH\'ES 64 (1986), pp.~53--110.
\begin{itemize}
\item The page range is confirmed by the scan.
\item \S3, Proposition, with the torus commutator formula credited to
[KP, Prop.~0.1.5].
\item \S9, Lemma and Proposition (good elements of the form $x^n$).
\item The Introduction and the Corollary in \S12 for $(n,N)=1$.
\end{itemize}
\item \emph{Zou, arXiv:2209.05852v3 (2024).}
\begin{itemize}
\item (2.2) (center);
\item \S2.3 ($d_r$);
\item (3.2) and Remarks~3.1 and~3.8 (the corrected scalar);
\item Proposition~3.6 (built on [FK, Theorem~26.1]);
\item the standing assumption $\gcd(n,p)=1$ in \S4.
\end{itemize}
Part~IV does not cite this paper at all. It should, because the relation
of $L_D$ to the GL-side correspondence goes through Zou's (3.2) (compare
Part~III, Remark~4.4).
\end{itemize}

\paragraph{Not verified by me (sources not attached).}
\begin{itemize}
\item \emph{[ZouTypes]} (J.\ Inst.\ Math.\ Jussieu 24 (2025)): the
bibliographic data, and the claim that it assumes $(p,n)=1$.
\item \emph{[Takeda]} (Pacific J.\ Math.\ 290 (2017)): the data, and what
``the scalar-center formula'' refers to.
\item \emph{[Kottwitz84], \S\S5, 7--9:} the Part~III log reports these
checked.
\item \emph{[Li14] = [Li4]:} the Part~II log reports it checked on EuDML.
I have not checked the DOI.
\item \emph{[DKV].}
\end{itemize}
[Cohen] and [Adams] are in the bibliography but are never cited.

\section*{4. Readability}
\begin{enumerate}[label=(R\arabic*)]
\item \emph{Drafting history must go to the log.} Three passages are about
the drafting process, not the mathematics:
\begin{itemize}
\item the end of the ``status labels'' paragraph (``Claude's proposed
cross-part review is therefore a particularly valuable prerequisite
\ldots'');
\item the acknowledgments' ``whom the author plans to consult separately
\ldots'';
\item the acknowledgments' ``should be checked against the exact cocycle
conventions of Part~I before publication''.
\end{itemize}
After \S1.2(d) above, the last of these is also resolved.
\item \emph{Disclaimers.} ``No \ldots\ is claimed'' appears in the
abstract, \S1, the remarks after Theorem~3.2, \S4, \S5, the ``preferred
order of work'' paragraph and the acknowledgments. State the status once,
in the table, as Parts~I and~II do.
\item \emph{Orienting paragraphs.} These are mostly present. \S2.1, \S3.1 and
\S4.1 begin abruptly.
\item \emph{Glossary.} As I said for Part~III: entries such as Hilbert symbol,
central character, induced torus and Tamagawa measure read as padding in
a research paper. Keep only a short notation list of terms specific to
the paper.
\item \emph{Abstract.} ``We compute the commutator radical \ldots\ and obtain
an explicit good-element criterion'' should become a citation of [FK, \S3].
\item \emph{One running example.} The two numerical remarks for
$(2,3,\mathbb Q_3(\zeta_3))$ should be merged and should cite Part~I,
Example~3.7, which uses the same data.
\end{enumerate}

\section*{5. Answers to the log's questions, and recommendations}
\begin{itemize}
\item[(Q1)] Yes: \S1.2(d).
\item[(Q2)] Yes, and the constants do not depend on the measures: \S1.5.
\item[(Q3)] Yes: \S1.1.
\item[(Q4)] Part~II's matching relates $\GDb$ to $\Gb$. Part~III's
pairing relates $D^\times$ to $\GDb$. They are different transfers, so no
compatibility between them is needed. When the two are composed, use
Part~II's measures on the $\GL$ side and any common measure on $H/Z$
(\S1.5).
\item[(Q5)] Yes: \S1.4.
\item[(Q6)] Yes: [FK, \S3, Proposition] eliminates the nonsplit-torus
task entirely (\S1.3).
\end{itemize}

\paragraph{Suggested Draft~2.}
\begin{enumerate}[label=\arabic*.]
\item An introduction that accurately states what Parts~I--III prove.
\item The wild problem as Part~II, goal~(ii), at wild places, via Zou's
(3.2) and your consistency report. Proposition~2.1 is reduced to a
remark citing Part~III, Corollary~5.2.
\item Good elements:
\begin{itemize}
\item [FK, \S3] for all tori, labeled \textbf{[E]};
\item the corollary of \S1.4, labeled \textbf{[P]}.
\end{itemize}
\item A precise statement of the $\GDb$--$D^\times$ comparison:
\begin{itemize}
\item the geometric sides, from Part~I, Proposition~6.1;
\item the local identities required at each place;
\item the product of local constants, including the kernels
$\mu_n^{k_v}/\mu_n$ at split places;
\item $(n,N)=1$.
\end{itemize}
Proving this would be a real theorem. Even a complete statement of the
reduction would be useful.
\item The explicit test for $(2,3,\mathbb Q_3(\zeta_3))$, carried out.
\begin{itemize}
\item Here $c_D=27$. With $\rho$ trivial, $L_D\Theta_\rho$ is
$3R(t)\tilde\theta(z)$ on $zt^*$ and~$0$ elsewhere.
\item Computing its coefficients needs a finite quotient
$\tH/\tZ U_k$ and its genuine character table.
\item This is a feasible GAP or Sage computation, and it decides the
first case of the wild problem.
\end{itemize}
\end{enumerate}

\end{document}
